\subsection{The configuration statements in the cited literature}

\begin{claim}{A.LIT1}{Proved: the literal scope of the disputed assertion}{A10, A14}
Mochizuki's original author Proposition~1.8(ii), equivalently
Proposition~1.6(ii) in the separate LaTeX version, asserts the full
prime-to-$p$ punctured-fibre identification when the ground field has
characteristic $p$ and $\Sigma$ consists of all primes different from $p$.
Consequently the noninjectivity established in the configuration audit
contradicts this specialization of the assertion itself. The February
2014 Comments do not amend this assertion.
\end{claim}
\begin{proof}
The conventions preceding original Proposition~1.6 define the geometric
configuration group by taking the maximal pro-$\Sigma$ quotient of the
tame fundamental group and then the maximal cuspidally pro-$\Sigma^\dagger$
quotient over the proper product. For the present finite-field choice,
$\Sigma^\dagger=\Sigma$. A closed subgroup of a pro-$\Sigma$ group is
already pro-$\Sigma$, so the latter quotient makes no change. The same
observation applies to the punctured fibre. Original Proposition~1.8(ii)
then identifies the inverse image of a rational-point decomposition
group under a configuration projection with that punctured group.
It is therefore not a statement about a singleton prime or about only
the completion of the fibre kernel.

The proof cites Stix's monodromy paper. Its relative completion retains
the entire base fundamental group and completes the fibre kernel.
The arXiv version's \S2.3 and Proposition~2.7 give that construction.
They do not justify additionally completing the base. The four corrections
in the 2014 Comments concern descended forms in decomposition recognition,
targets of diagonal covers, a Kummer injection, and infinitude of a
cyclotomic image, respectively. None replaces the fibre identification.
These readings establish the precise conflict; they do not establish a
verdict about every theorem of the paper.
\citep[original pp.~11,14--15]{LA_Mochizuki2007}
\citep{LA_Mochizuki2014,LA_Stix2005}
\end{proof}

\paragraph{Version control.}
The following numbering distinctions matter. In Sa\"idi--Tamagawa 2009,
the supplied arXiv Remark~3.13 is journal Remark~3.14, and the supplied
arXiv Theorem~2.7 is journal Theorem~2.6. The former is the explicitly
conditional configuration-space route to elimination of exceptional
points; journal Remark~3.13 concerns a different extension of the Isom
theorem. Mochizuki's original Theorem~3.10 is Theorem~3.1 in the LaTeX
version. Statements below specify the original numbering when necessary.
\citep{LA_ST2009,LA_Mochizuki2007}

\begin{claim}{A.LIT2}{Proved: a limited repair of the fibre input}{A10}
Let $k$ be a perfect field, let $C/k$ be a proper smooth
geometrically connected curve of genus at least two, and let
$\Sigma^\dagger=\{r\in\Sigma:r\text{ is invertible in }k\}$
be nonempty. Thus $\Sigma^\dagger=\Sigma$ in characteristic
zero and $\Sigma^\dagger=\Sigma\setminus\{p\}$ in
characteristic $p>0$.
Let $F\to E\to\Delta_2\to1$ be the geometric right-exact configuration
sequence, let $F\twoheadrightarrow\Delta_1$ be the proper fibre quotient,
and let $E\twoheadrightarrow\Delta_1\times\Delta_2$ be the proper product
quotient. Write $K_F,K_E$ for their cuspidal kernels. Suppose the relative
groups are formed from the \emph{tame} boundary groups, with Mochizuki's
maximal cuspidally pro-$\Sigma^\dagger$ convention. The proper groups
may retain $p$ in $\Sigma$. The relative cuspidally pro-$r$ fibre
theorem of Part~A, including \statementref{A.CF12}, holds for every
$r\in\Sigma^\dagger$.
Then
\[
 K_F^{(r)}\xrightarrow{\sim}K_E^{(r)}\quad(r\in\Sigma^\dagger),
 \qquad K_F^{\mathrm{ab}}\xrightarrow{\sim}K_E^{\mathrm{ab}}.
\]
Consequently the configuration fibre sequence is left exact after
cuspidal abelianization, as well as after the specified fixed-prime
cuspidal completion. The one-puncture geometrically cuspidally central
fibre identification also holds. For a $k$-rational point, the
corresponding arithmetic inverse-image identifications hold over
the unchanged quotient $G_k$.
\end{claim}
\begin{proof}
Right exactness gives a surjection $K_F\twoheadrightarrow K_E$:
an element of $K_E$ lifts to $F$, and its trivial first proper projection
forces the lift into $K_F$. The relative pro-$r$ theorem identifies the
two cuspidal pro-$r$ kernels by taking kernels of the identified fibre
groups over $\Delta_1$. Every abelian pro-$\Sigma$ group is the product
of its pro-$r$ factors. In positive characteristic there is no $p$
factor in the present cuspidal
kernels: each is topologically normally generated by conjugates of
tame inertia, and every map to a finite $p$-group kills each such
generator. In particular $K_F^{(p)}=K_E^{(p)}=1$ when $p\in\Sigma$.
This assertion would be false for the full \emph{wild} punctured
\'etale group, and that group is not the convention used here.
Moreover the maximal abelian pro-$r$ quotient of
a group is the abelianization of its maximal pro-$r$ quotient. Taking
the product of the resulting isomorphisms proves the assertion about
abelianizations. Applying the elementary extension isomorphism criterion
over $\Delta_1$ proves the cuspidally abelian fibre statement.

For one puncture, taking $\Delta_1$-coinvariants of the abelian cuspidal
kernel leaves the orientation module
$\widehat{\Z}^{\Sigma^\dagger}(1)$, generated by the boundary loop. The other
geometric proper factor acts trivially on this orientation: the family
has one labelled boundary and its local parameter gives the same tame
inertia orientation. Hence taking its coinvariants adds no relation.
This proves the geometrically central assertion. The arithmetic action
on the surviving module is the cyclotomic action. Equivalently this
last extension is constructed directly from the Kummer cycle class of
the diagonal, whose restriction to the fibre is the degree-one point
class. Neither argument requires injection of the full fibre.

The relative pro-$r$ proof is geometric and uses no finite-field
count or Frobenius argument. In positive characteristic it is
exactly the direct Stix argument of \statementref{A.CF12}; in
characteristic zero the same argument applies to every $r$, or
one uses the topological construction after descending the finite
data. For a rational point, the natural map from the punctured
arithmetic quotient to the inverse image of its decomposition
group is an isomorphism on the geometric kernels just proved and
on the arithmetic quotients, both $G_k$. The elementary extension
criterion makes it an isomorphism. Finite sets of closed points
are handled after a finite separable splitting extension; all
maps are natural geometric maps over the unchanged arithmetic
quotient, so their Galois descent data commute and the
identification descends. This justifies the stated field scope.
\end{proof}

\paragraph{Exact repair boundaries.}
This proves a replacement for the false intermediate assertion over
every perfect ground field with these tame cuspidal conventions,
including proper $\Sigma$-groups which retain the positive
characteristic. It covers all fields of the original paper:
its \S0 defines a nonarchimedean local field to mean a finite
extension of $\mathbb Q_p$, so its local convention
$\Sigma^\dagger=\Sigma$ is exactly the characteristic-zero
case here. It does not turn the original
full sequence into an exact sequence. The following dependency loci were
checked in the original author text. Proposition~1.8(iii),(iv) uses
1.8(ii) for central quotients. The paragraph preceding Proposition~1.12
(original p.~20) cites 1.8(ii),(iii) for surjectivity of the symmetrized
fundamental section. Proposition~1.14 builds cuspidally abelian quotients
as inverse limits. Before Proposition~3.8 (original p.~63), a punctured
fibre is inserted in the configuration group by 1.8(ii); Proposition~3.9(v)
identifies the fixed-$l$ cuspidal quotient, which enters Theorem~3.10.

At these loci one may use the central or abelian fibre identification
when the subsequent construction is central or abelian, and the relative
pro-$l$ identification when it is cuspidally pro-$l$. This specifies the
replacement actually proved here. Constructions on finite proper covers
retain their labelled boundary components. Here is the needed extension
of the relative proof. For a product compactification $C_0\times B$
with finitely many disjoint section graphs, choose the finite Galois
$A$-cover $C_1\to C_0$ required by a finite fibre quotient. Take a
finite Galois $\Sigma$-cover of $B$ dominating its pullbacks along
every section; their finite intersection and core give such a cover.
All lifted punctures are then labelled, the compactification remains
a product, and mod-$r$ monodromy is a shear between proper homology
and the boundary module. The same relative pro-$r$ and normal-core
argument applies, with finite proper deck group $A\times A_B$.
When this group's order is divisible by $p$, the direct algebraic
argument of \statementref{A.CF12} uses Stix's relative pro-$r$
completion over the full base; it requires no equivariant lift of the
proper cover to characteristic zero.
This proves the required compatibility for these finite proper covers.
For naturality under a finite proper pullback not itself visible
as a $\Sigma$-cover, apply the construction to that curve with
its own proper pro-$\Sigma$ group: naturality of the restriction
maps and the diagonal cycle class gives the comparison square.
One does not pretend that this cover is an open subgroup of the
original pro-$\Sigma$ group. Any conclusion retaining the entire
mixed-prime punctured kernel still loses its asserted fibre
identification. The dependency census below records that loss
separately from the repaired quotient constructions.
\citep[Propositions~1.8,1.12,1.14,3.9 and Theorem~3.10]{LA_Mochizuki2007}

\begin{claim}{A.LIT3}{Proved: the primary and total formation distinction}{A10}
Mochizuki--Tamagawa's configuration exactness theorem and Wakabayashi's
configuration theorem do not assert the mixed prime-to-$p$ exactness
refuted in Part~A.
\end{claim}
\begin{proof}
Mochizuki--Tamagawa Definition~1.1 restricts the formation to a
PT-formation, namely primary or total. Proposition~2.2 additionally
requires invertibility on the connected base and the stated formation
conditions on the cyclotomic and compactified-fibre mod-$l$ monodromy
images. For classes of all finite $\Sigma$-groups, primary means one
prime; total means all finite groups. The proof uses that the kernel of
the automorphism action of a finitely generated pro-$l$ group on its
mod-$l$ abelianization is pro-$l$. A set of several primes omitting the
characteristic is not one of these two formations. The 2015 Comments
do not enlarge the formation scope. Wakabayashi fixes one invertible
prime throughout. Thus neither theorem supplies the disputed
mixed-prime completion step.
\citep[Definition~1.1 and Proposition~2.2]{LA_MT2008}
\citep{LA_MT2015,LA_Wakabayashi2016}
\end{proof}

\begin{claim}{A.LIT4}{Proved: Wakabayashi's actual input and output}{A11}
Fix a prime $l$ invertible in the two finite fields. Wakabayashi's
Theorem~A, proved in the more general logarithmic form as Theorem~3.9,
starts with a Frobenius-preserving isomorphism of the geometrically
pro-$l$ groups of once-punctured curves, carrying the specified new
cusp decomposition groups to one another. It extends this isomorphism
to the two-point configuration groups, compatibly with the proper
quotient and switching, up to the stated inner ambiguity. Corollary~3.10
transports punctures at other points with corresponding proper
decomposition groups. Theorem~C, namely Theorem~4.4, reconstructs the
higher configuration groups from the geometrically pro-$l$ curve group.
These are reconstruction results in the singleton-prime category.
\end{claim}
\begin{proof}
These are the hypotheses and conclusions of the supplied source. Its
proof uses the inverse image of a point decomposition group under the
configuration projection, so the fibre method itself is already present.
Remark~3.10.1 explicitly permits nonuniqueness of the corresponding
point. The theorem retains a pro-$l$ proper group, rather than the full
prime-to-$p$ proper group with only a pro-$r$ cuspidal kernel. No operation
in the theorem reconstructs mixed finite quotients from separate
pro-$l$ quotients. Nor may the cofinite-prime exceptional-set theorem
be applied to a singleton set without verifying its hypothesis. Thus
these results do not conflict with AG.40--AG.42 and, with the stated
Sa\"idi--Tamagawa finiteness input alone, do not supply their universal
point-separation conclusion.
\citep[Theorems~3.9,4.4 and Corollary~3.10]{LA_Wakabayashi2016}
\end{proof}

\paragraph{Sa\"idi--Tamagawa dependencies.}
The 2009 configuration route in journal Remark~3.14 is expressly
conditional. Its failure does not by itself invalidate that paper's
Isom theorem. Theorem~2.2 reconstructs the configuration cuspidally
abelian group; Corollary~2.3 restricts the labelled puncture comparison
to the complements of exceptional sets; journal Theorem~2.6 gives
fixed-$l$ cuspidal reconstruction. The 2017 text repeats these as
Theorem~3.2, Corollary~3.3 and Theorem~3.6. Thus their imported
configuration machinery depends on the central, abelian, and fixed-$l$
consequences just distinguished, not on a final theorem asserting
full mixed-prime configuration exactness.

There is a further scope distinction. Original Mochizuki Theorem~3.10
allows arbitrary matching rational points, without a separation
hypothesis. Its proof begins with configuration cuspidally abelian
reconstruction from Theorem~1.16(iii), chooses matching splittings in
the diagonal decomposition groups, and restricts to the corresponding
cuspidally abelian fibres. Thus the reference to an outside-exceptional-set
corollary in Sa\"idi--Tamagawa is not on its own proof that the original
arbitrary-point theorem has no such argument. The needed repair is the
fibre input at the stated quotient level. Our main proof uses its own
relative fibre and transport constructions and does not assume the
published reconstruction theorem in order to prove that repair.
\citep{LA_ST2009,LA_ST2017,LA_Mochizuki2007}

\subsection{Proper support fields and the later exceptional-set theorem}

\begin{claim}{A.LIT5}{Proved: the exact later largeness conclusion}{A10, A14}
In Sa\"idi--Tamagawa 2017, Definition~2.6 calls $\Sigma$ $J_C$-large
when the omitted-prime representation
$\rho_{J_C,\Sigma'}:G_k\to\Aut(T_{\Sigma'}J_C)$ is noninjective.
Definition~2.9 calls $C$ almost $\Sigma$-separated when its exceptional
closed points are supported over some proper algebraic subfield
$F\subsetneq\bar k$. Proposition~2.12 proves this latter conclusion
under $J_C$-largeness. It does not in general assert that the
exceptional set is finite.
\end{claim}
\begin{proof}
Put $F_0=\bar k^{\ker\rho_{J_C,\Sigma'}}$. The proof uses the difference
map
\[
 \delta:C\times C\longrightarrow J_C,\qquad (P,Q)\longmapsto[P-Q].
\]
Off the diagonal its geometric fibres have cardinality at most two,
or at most one when $C$ is not hyperelliptic. Indeed, two different
pairs in a fibre give $P+Q'\sim P'+Q$. A nonconstant degree-two linear
system is the hyperelliptic system; in that case the only possible
second pair is $(\iota Q,\iota P)$, with $\iota$ the hyperelliptic
involution. This argument is the predecessor explicitly cited from
2009 Claim~1.9(i), where the curve has first been made
nonhyperelliptic.

An exceptional pair has distinct images on $C_{\bar k}$ and its
difference lies in $J_C(\bar k)\{\Sigma'\}\subset J_C(F_0)$.
Taking the subgroup of squares of $G_{F_0}$ kills the possible
two-element permutation. Hence both coordinates are defined over a
finite extension of $F_0$, still proper in $\bar k$. This is exactly
the argument of Proposition~2.12. Proposition~2.13 then chooses a prime
$l$ for which the corresponding $\Z_l$ tower has finite intersection
with every fixed finite extension of that support field. It may choose
$l\in\Sigma\cup\{p\}$ under largeness, and $l\in\Sigma$ if the
Jacobian also has positive $p$-rank. Proposition~2.11(i) gives
$E_{C'}=f^{-1}(E_C)$ for finite \'etale $\Sigma$-covers and transports
almost separation through them.
\citep[Definitions~2.6,2.9 and Propositions~2.11--2.13]{LA_ST2017}
\end{proof}

\begin{claim}{A.LIT6}{Proved with the relative criterion of Part A}{A10, A14}
The separation argument extends from finite exceptional support to
almost $\Sigma$-separation whenever $\Sigma\setminus\{p\}$ is nonempty,
including $p\in\Sigma$. We use the relative fixed-$r$ fibre and
transport statements of Part~A in that $\Sigma$-category, as extended
by \statementref{A.CF12}. Consequently the $J_C$-large
almost-separated cases of Sa\"idi--Tamagawa become separated.
\end{claim}
\begin{proof}
Fix $r\in\Sigma\setminus\{p\}$. The point-transport and difference-cocycle argument
in Part~A gives
$\kappa_r(f(R)/a)=0$ outside the exceptional support and the two
punctures, where $a\in k^\times$ and
$\operatorname{div}(f)=n(P-Q)$. Here one may take $n$ to be the
\emph{least} positive order of $[P-Q]$ in $J_C(\bar k)$; it need not
be invertible at $r$. The expression with zero residues is
$\kappa_r(f)-n\chi$, so no division by $n$ occurs. Hilbert~90 descends
the generator of this rational principal divisor to $k(C)$.

The function $f$ is not an $r$th power over $\bar k(C)$. If $r\nmid n$,
its valuations prove this; if $r\mid n$, an $r$th root would make
$(n/r)(P-Q)$ principal, contrary to minimality. Choose a finite field
$L\supset k$ containing $\mu_r$ and let $D/L$ be the smooth proper
model of $z^r=f/a$. This curve is geometrically connected. Let $h$ be
its genus and let $L_m/L$ run through the powers of the auxiliary
prime $l$ supplied by Proposition~2.13. The exceptional points rational
over these fields lie in a fixed finite set.

Write $Q_m=|L_m|$ and let $M_m$ count points of
$C(L_m)\setminus\{P,Q\}$ at which $f/a$ is an $r$th power.
Every such point has $r$ rational inverse images on $D$, so
$rM_m\leq|D(L_m)|$. The Weil bounds imply
\[
 |C(L_m)|-2-M_m\ \geq\
 (1-r^{-1})(Q_m+1)-(2g+2h/r)\sqrt{Q_m}-2,
\]
which tends to infinity. Some point outside the bounded exceptional
set therefore has nonzero mod-$r$ Kummer class, contradicting the
displayed evaluation. This proves the extension. The counting argument
does not require $l\ne r$.
\end{proof}

\paragraph{The elementary full prime-to-$p$ specialization.}
For the main theorem one can avoid the earlier Boxall-based finiteness
input entirely. Let $f_p$ be the $p$-rank of $J_C$. Choose an odd prime
$l\ne p$ not dividing $|\mathrm{GL}_{f_p}(\F_p)|$, with the evident
convention when $f_p=0$. The subgroup $\Z_l\subset G_k$ acts trivially
on $T_pJ_C$, since the congruence kernel in its automorphism group is
pro-$p$, and therefore acts trivially on all geometric $p$-primary
torsion. The preceding two-point difference-fibre bound then puts
every distinct pair with $p$-primary difference over
$\bar k^{\Z_l}$. Along the $l$-power constant-field tower these bad
labels are uniformly bounded. This is a direct special case of the
known support-field method, combined with the new criterion as proved
in Part~A; the difference-fibre lemma itself is not new.

\subsection{What is and is not supplied by the opened sources}

\begin{claim}{A.LIT7}{Proved: the recorded predecessor comparison}{A14}
The following comparison records exact predecessors and the additional
work required in this audit. It is not a priority claim.
\end{claim}
\begin{center}\small
\begin{longtable}{p{0.17\linewidth}p{0.37\linewidth}p{0.38\linewidth}}
\toprule
Step&Opened predecessor&Additional issue resolved or retained\\\midrule
AG.40, point recognition&
Mochizuki original Remark~1.18.2 and its 2014 correction; the conclusion
is ``Any Frobenius-preserving isomorphism $\alpha$ is quasi-point-theoretic.''&
One must prove that the relative quotient still recovers the compactified
$r$-adic cohomology of every visible cover; the finite-intersection
argument must allow all relevant descended forms.\\
AG.40, Kummer evaluation&
Mochizuki Proposition~2.1(ii),(iii): residue exactness and evaluation of
a function's Kummer class at a point. Sa\"idi--Tamagawa repeats it.&
The over-identity cusp isomorphism yields a difference cocycle with
residues $+1,-1$; combining it with transported point labels gives the
contradiction. No identical full criterion was located.\\
AG.41, exactness&
Stix's relative completion; M--T Proposition~2.2 in PT formations;
Hoshi--Sawada--Tsujimura Definition~1.6(ii) and Corollary~1.14;&
Neither states the mixed proper-base, fixed-$r$ cuspidal quotient used
here. The direct proof and its finite-quotient checks remain essential.\\
AG.42, fibre transport&
Mochizuki original Remark~1.7.1 at the central level; Wakabayashi
Corollary~3.10 for a fixed prime; Sun Definition~3.1 and Corollary~3.8.&
The general pullback trick is known. Its use in the relative mixed-base
quotient also requires the local diagonal-boundary lift and compatibility.\\
Specialization and inertia&
SGA1 XIII Corollary~2.9, \S2.10, Lemma~2.11 and Proposition~5.2.&
The hypotheses require a smooth proper pair with relative normal crossings
and invertible coefficient primes; the local Kummer extension exists on
the whole regular chart, by purity.\\
Orbifold completion&
Boggi--Zalesskii Definition~1.1 and Lemma~3.1, with a torsion-free normal
subgroup and quotient in the extension-closed formation.&
The prime-to-$p$ formation meets these hypotheses when the chosen
orbifold quotient has prime-to-$p$ order; no $p$-acyclicity is required.\\
\bottomrule
\end{longtable}
\end{center}
The 2024 Hoshi--Sawada--Tsujimura result requires good log-smooth
pairs and a proper geometrically connected morphism, smooth away
from the base boundary, with smooth relative boundary there, and
log smooth as a morphism of pairs. Definition~1.6(ii) takes the
prime-to-characteristic quotient of the \emph{kernel} of the map
between full tame fundamental groups. Corollary~1.14 identifies
the geometric prime-to-characteristic fibre when it is centre-free.
The base remains the full tame group. This is a closer recent
predecessor for relative completion, and supplies no simultaneous
prime-to-$p$ completion of that base. Its proof in \S1 was checked
at these precise statements.
\citep[Definitions~1.1,1.3,1.6 and Corollary~1.14]{LA_HST2024}

The October 2026 Hoshi--Tsujimura preprint also has a point-section
injectivity theorem: Theorem~A(i), proved as Theorem~1.8, concerns
affine varieties over a characteristic-$p$ field admitting a Galois
extension of degree $p$. It uses the full \'etale fundamental group
and Artin--Schreier evaluations. It does not assert the proper-curve
prime-to-$p$ injectivity at issue here.
\citep[Theorem~A(i) and Theorem~1.8]{LA_HT2026}

The sources were opened in full or at the identified proofs, not merely
through search snippets. Searches included each named paper, both Sun
preprints, and combinations of prime-to-$p$ point separation, exceptional
sets, cuspidalization, and configuration exactness. No independently
published universal prime-to-$p$ point-separation proof was found in
that bounded search. This does not prove novelty or exclude an unknown
predecessor.\citep{LA_Mochizuki2007,LA_Mochizuki2014,LA_ST2009,
LA_ST2017,LA_MT2008,LA_Stix2005,LA_Wakabayashi2016,LA_Sun2025,
LA_SGA1,LA_BoggiZalesskii2019,LA_HST2024,LA_HT2026}

\subsection{A finite dependency census for the disputed exactness}

\begin{claim}{A.LIT8}{Proved}{Which imported fibre arguments require replacement}
For the five papers specified in A10, the dependencies on the
configuration fibre assertion are as listed below. At every
listed central, abelian, or fixed-prime use, A.LIT2 supplies the
same quotient statement with its maps and boundary data. The
literal mixed-prime fibre insertion and any argument requiring
that full insertion are lost. This is a census of the disputed
dependency, not a reproof of unrelated theorems in those papers.
\end{claim}
\begin{center}\small
\begin{longtable}{p{0.23\linewidth}p{0.35\linewidth}p{0.34\linewidth}}
\toprule
Source locus&Use of the fibre input&Verified replacement or exact loss\\\midrule
Mochizuki 2007, Proposition~1.8(ii)&
Identifies the full punctured group with a configuration inverse
image; concludes slimness of the full configuration group.&
The full mixed-prime identification is false in the specialization
of A.LIT1. Only its central, abelian and fixed-prime versions are
proved here. Full mixed-group slimness is not inferred from those
versions or from the invalid full-fibre proof.\\
Proposition~1.8(iii),(iv)&
Specializes the fundamental section, proves surjectivity from
boundary inertia, and identifies central quotients.&
The diagonal cycle class and its point restrictions give the
same maps. Primitive boundary inertia gives surjectivity and
the universal central quotient, by A.LIT2. The separate slimness
claim for curve groups in (iii) is proved there from curve-group
results, not from the disputed configuration insertion.\\
Paragraph before Proposition~1.12; Propositions~1.12,1.14--1.16,
1.19&
Surjectivity of symmetrized fundamental sections; inverse limits
of central quotients; reconstruction and switching of the
cuspidally abelian configuration group.&
Apply the repaired central maps on each finite proper pullback
and take their compatible inverse limits. The needed comparison
and rigidity are at the abelian level.\\
Paragraph before Proposition~3.8; Propositions~3.8--3.9;
Theorem~3.10&
Writes the full punctured fibre as a subgroup to define diagonal
actions, then uses its fixed-$l$ Lie quotients and finite proper
pullbacks.&
Replace the ambient group and the fibre from the outset by their
relative cuspidally pro-$l$ quotients. A.LIT2 and the retained
boundary give those actions and inclusions. Every Lie object
used here factors through this quotient. The displayed full
mixed-group inclusion on original p.~63 remains invalid.\\
Mochizuki--Tamagawa 2008, Proposition~2.2 and its subsequent
configuration sequences&
Uses its own PT-formation homotopy theorem. The citation to
Mochizuki 2007 in Remark~4.8.2 concerns a proposed later
application.&
No proof in this dependency chain invokes the false mixed
formation. Its singleton and total hypotheses are retained.\\
Sa\"idi--Tamagawa 2009, Theorem~2.2, Corollary~2.3,
Proposition~2.4, Theorem~2.6&
Imports Mochizuki's abelian reconstruction, marked restrictions,
central quotient, and fixed-$l$ reconstruction, respectively.&
These are exactly the repaired quotient inputs. The conditional
full-group proposal in journal Remark~3.14 is lost; it is not a
premise of their Isom proof.\\
Sa\"idi--Tamagawa 2017, Theorem~3.2, Corollary~3.3,
Proposition~3.4, Theorem~3.6&
Reuses the preceding four constructions, with its own exceptional
sets and largeness hypotheses.&
The same four quotient repairs apply. Its support-field and
function-reconstruction arguments have independent hypotheses,
which remain as in C.LIT5--7.\\
Wakabayashi 2016, Proposition~1.4; Theorems~3.9,4.4&
The fibre theorem comes from M--T Proposition~2.2 or Stix in
the single-prime category. The proper two-point case of 4.4
cites Mochizuki's fixed-$l$ theorem.&
No full mixed-prime fibre assertion is needed in the category
used by these proofs. Its imported fixed-$l$ theorem is in the
repaired scope, and also in the singleton case where the
homotopy theorem is valid.\\
\bottomrule
\end{longtable}
\end{center}
\begin{proof}
The full-text occurrence census for original Proposition~1.8(ii)
has three direct loci: its use in the proof of (iii), the
surjectivity paragraph before Proposition~1.12, and the fibre
insertion immediately before Proposition~3.8. The other explicit
citations to Proposition~1.8 use its proper-group slimness in
(i), its curve or central assertions in (iii),(iv), or the
definitions of the punctured groups. We checked the mathematical
content of these loci, not just the citation strings.

For the first two loci, the source's required map sends each
boundary generator to its corresponding orientation-module
generator in the extension defined by the diagonal class.
A.LIT2 constructs this map without a full-fibre injection. Since
the central cuspidal kernel is generated by these images, the
map is surjective and has the universal central quotient kernel.
On a finite proper cover, conjugate boundary components give the
separate factors of the symmetrized extension. The labelled-cover
argument in A.LIT2 proves compatibility with their deck actions.
Taking inverse limits identifies the full cuspidal abelian
module, hence gives the identifications in Proposition~1.14.
The proof of Theorem~1.16(iii) reconstructs these extensions and
their transition maps, and its uniqueness argument uses
Proposition~1.15 at this abelian level. The induced geometric
switching map exists before taking any quotient; its needed
abstract characterization in Proposition~1.19 likewise uses
Proposition~1.15, not a full mixed-group fibre theorem.

At the third locus the paper has $S=S_* =\{x_*\}$, hence
$S_0=\varnothing$. Its ensuing objects are constructed from the
pro-$l$ groups of the punctured finite proper covers, their
weight quotients and their Lie realizations. Consequently each
action in Propositions~3.8--3.9 factors through the relative
cuspidally pro-$l$ fibre with its proper pro-$\Sigma$ quotient.
Replace the displayed full configuration group by that relative
quotient. A.LIT2 identifies its fibre with the relative
punctured group and identifies the image of the diagonal
decomposition group with its retained tame boundary. Conjugation
therefore supplies the same continuous diagonal actions on every
pro-$l$ object. Naturality on finite proper pullbacks supplies
their transition maps. No weight filtration, trace map or
Frobenius splitting in the subsequent proof changes. The
slimness used at the end of Proposition~3.9(v) is explicitly
slimness of the \emph{singleton-$l$} curve and configuration
groups, where M--T Proposition~2.2 applies. Thus the fixed-$l$
reconstruction of Theorem~3.10 has the repaired input, without
restoring the original full mixed-group insertion.

The two Sa\"idi--Tamagawa papers cite precisely these central,
abelian and fixed-$l$ results at the four listed theorem loci.
Their subsequent use in Green compatibility and multiplicative
function reconstruction only uses these quotients; C.LIT5--7
give the separate consequences of removing exceptional sets.
For M--T, the full text restricts configuration exactness to
PT formations in Proposition~2.2. Its sole invocation of the
cuspidalization paper in Remark~4.8.2 is motivational, not a
proof input. Finally Wakabayashi fixes one invertible prime
throughout and gives the exact source of its fibre comparison
in Proposition~1.4; its references to Mochizuki's weight
filtration and fixed-$l$ reconstruction stay in that category.
These checks prove the stated dependency classification.
\citep[Propositions~1.8,1.12,1.14--1.16,1.19,3.8--3.9;
Theorem~3.10]{LA_Mochizuki2007}
\citep[Proposition~2.2, Remark~4.8.2]{LA_MT2008}
\citep[\S2, Remark~3.14]{LA_ST2009}
\citep[\S3]{LA_ST2017}
\citep[Proposition~1.4, Theorems~3.9,4.4]{LA_Wakabayashi2016}
\end{proof}
